\section{Spectral cut-offs in Hastings' approximate ground-state projector}
\label{sec:hastings_spectral_cutoff}

% Hastings, arXiv:0705.2024v4, law16.tex, lines 829--841 (step four of the
% proof of Lemma 1, equation obd).

In step four of the proof of \cite[Lemma~1]{Hastings2007AreaLaw}, the localized
filtered operators $M_L$ and $M_R$ nearly annihilate $\Psi_0$, and $O_L,O_R$
are their spectral projections onto $[-\theta,\theta]$ (two-sided rather than
$(-\infty,\theta]$ because $M_L,M_R$ need not be positive and the proof bounds
$\|M_LO_L\|$~\cite{gap:hastings07_spectral_cutoff_threshold}).

\begin{definition}[Spectral cut-off]
    \label{def:hastings_spectral_cutoff}
    \lean{TNLean.MPS.AreaLaw.Hastings.spectralCutoff}
    \leanok
    For $M=M^\dagger$ and $\theta\in\R$, the spectral cut-off is
    $O_\theta(M)=\mathbf 1_{[-\theta,\theta]}(M)$.
\end{definition}

\begin{theorem}[The cut-off is an orthogonal projection]
    \label{thm:hastings_spectral_cutoff_projection}
    \lean{TNLean.MPS.AreaLaw.Hastings.isStarProjection_spectralCutoff}
    \leanok
    \uses{def:hastings_spectral_cutoff}
    For $M=M^\dagger$ and $\theta\in\R$, $O_\theta(M)^2=O_\theta(M)=O_\theta(M)^\dagger$.
\end{theorem}
\begin{proof}\leanok
    The indicator $f=\mathbf 1_{[-\theta,\theta]}$ is real-valued with $f^2=f$.
\end{proof}

\begin{theorem}[A cut-off moves a nearly annihilated vector little]
    \label{thm:hastings_spectral_cutoff_obd}
    \lean{TNLean.MPS.AreaLaw.Hastings.norm_sub_spectralCutoff_le}
    \leanok
    \uses{def:hastings_spectral_cutoff}
    For $M=M^\dagger$, $\theta>0$, and any vector $\Psi$,
    \begin{align}
        \lVert\Psi-O_\theta(M)\Psi\rVert
         & \le\frac{\lVert M\Psi\rVert}{\theta}.
        \label{eq:hastings_spectral_cutoff_obd}
    \end{align}
    For $\lVert M_L\Psi_0\rVert\le C\Delta E^{-1}J^2e^{-l/(3\xi)}$ and
    $\theta=\Delta E^{-1}J^2e^{-l/(6\xi)}$, (\ref{eq:hastings_spectral_cutoff_obd})
    gives $\lVert(O_L-1)\Psi_0\rVert\le Ce^{-l/(6\xi)}$ as in
    \cite[eq.~(obd)]{Hastings2007AreaLaw}, and similarly for $O_R$.
\end{theorem}
\begin{proof}\leanok
    Set $h(t)=t^{-1}\mathbf 1_{\lvert t\rvert>\theta}(t)$, so
    $1-\mathbf 1_{[-\theta,\theta]}(t)=t\,h(t)$ and $\lvert h(t)\rvert\le\theta^{-1}$.
    Then $1-O_\theta(M)=h(M)M$ with $\lVert h(M)\rVert\le\theta^{-1}$, giving
    (\ref{eq:hastings_spectral_cutoff_obd}).
\end{proof}

\begin{theorem}[The cut-off reduces the operator]
    \label{thm:hastings_spectral_cutoff_norm}
    \lean{TNLean.MPS.AreaLaw.Hastings.norm_mul_spectralCutoff_le}
    \leanok
    \uses{def:hastings_spectral_cutoff}
    For $M=M^\dagger$ and $\theta\ge0$, $\lVert M\,O_\theta(M)\rVert\le\theta$.
\end{theorem}
\begin{proof}\leanok
    Write $M\,O_\theta(M)=g(M)$ with $g(t)=t\,\mathbf 1_{[-\theta,\theta]}(t)$
    and $\lvert g(t)\rvert\le\theta$.
\end{proof}

\begin{theorem}[The cut-off estimate]
    \label{thm:hastings_spectral_cutoff_pair_norm}
    \lean{TNLean.MPS.AreaLaw.Hastings.norm_add_mul_mul_le_of_isStarProjection,
        TNLean.MPS.AreaLaw.Hastings.norm_add_mul_mul_spectralCutoff_le}
    \leanok
    \uses{thm:hastings_spectral_cutoff_projection,thm:hastings_spectral_cutoff_norm}
    If $O,O'$ are commuting orthogonal projections with $\lVert AO\rVert\le a$
    and $\lVert BO'\rVert\le b$, then $\lVert(A+B)OO'\rVert\le a+b$. In
    particular, for $M_L=M_L^\dagger$, $M_R=M_R^\dagger$, $\theta\ge0$, and
    commuting $O_L=O_\theta(M_L)$, $O_R=O_\theta(M_R)$,
    $\lVert(M_L+M_R)O_LO_R\rVert\le2\theta$.
\end{theorem}
\begin{proof}\leanok
    Use $(A+B)OO'=(AO)O'+(BO')O$, $\lVert O\rVert,\lVert O'\rVert\le1$, and
    Theorem~\ref{thm:hastings_spectral_cutoff_norm}.
\end{proof}

\begin{theorem}[Two cut-offs move a vector by at most the sum]
    \label{thm:hastings_spectral_cutoff_two}
    \lean{TNLean.MPS.AreaLaw.Hastings.norm_sub_mul_le_of_isStarProjection}
    \leanok
    For an orthogonal projection $O$, any operator $O'$, and any vector $\Psi$,
    $\lVert\Psi-OO'\Psi\rVert\le\lVert\Psi-O\Psi\rVert+\lVert\Psi-O'\Psi\rVert$.
\end{theorem}
\begin{proof}\leanok
    Write $\Psi-OO'\Psi=(\Psi-O\Psi)+O(\Psi-O'\Psi)$ and use $\lVert O\rVert\le1$.
\end{proof}

\begin{theorem}[Support of the cut-off]
    \label{thm:hastings_spectral_cutoff_support}
    \lean{TNLean.MPS.AreaLaw.Hastings.spectralCutoff_localLift,
        TNLean.MPS.AreaLaw.Hastings.isSupportedOn_spectralCutoff}
    \leanok
    \uses{def:hastings_spectral_cutoff}
    For $K=K^\dagger$ on a region $D$, $O_\theta(K\otimes1_{D^c})=O_\theta(K)\otimes1_{D^c}$;
    thus $O_\theta(M)$ is supported on $D$ whenever $M=M^\dagger$ is.
\end{theorem}
\begin{proof}\leanok
    The unital $*$-homomorphism $K\mapsto K\otimes1_{D^c}$ commutes with the
    functional calculus on the finite spectrum of $K$.
\end{proof}

\begin{theorem}[Cut-offs on disjoint regions commute]
    \label{thm:hastings_spectral_cutoff_commute}
    \lean{TNLean.MPS.AreaLaw.Hastings.commute_spectralCutoff_localLift_of_disjoint,
        TNLean.MPS.AreaLaw.Hastings.commute_spectralCutoff_of_isSupportedOn}
    \leanok
    \uses{thm:hastings_spectral_cutoff_support}
    If $M=M^\dagger$ and $M'=(M')^\dagger$ are supported on disjoint regions
    $X$ and $Y$, then $[O_\theta(M),O_{\theta'}(M')]=0$ for all
    $\theta,\theta'\in\R$; in particular, $[O_L,O_R]=0$.
\end{theorem}
\begin{proof}\leanok
    By Theorem~\ref{thm:hastings_spectral_cutoff_support}, $O_\theta(M)$ and
    $O_{\theta'}(M')$ are supported on disjoint regions $X$ and $Y$, hence
    commute.
\end{proof}
